\section*{Цель}
Получение практических навыков по тестированию клиентской части веб-приложения.

\section*{Задачи}
\begin{enumerate}
    \item Провести тестирование совместимости веб-системы
    \item Протестировать адаптивность веб-системы.
    \item Проверить валидность верстки.
    \item Протестировать структуру веб-приложения.
    \item Провести технический аудит веб-системы
\end{enumerate}

\section*{Описание веб-приложения}
Для выполнения работы был выбран сайт компании «КЗ Групп», доступный по адресу \url{https://kzgroup.ru/}.

\textit{Предметная область.} Поставка промышленного оборудования и комплектующих для различных предприятий.

\textit{Назначение веб-приложения.} Сайт является визитной карточкой компании и содержит каталог продукции. Он предназначен для того, чтобы пользователи могли:
\begin{itemize}
\item Посмотреть список доступного оборудования и его характеристики.
\item Узнать о предоставляемых услугах.
\item Найти контакты организации или оставить заявку на обратный звонок.
\end{itemize}

Сайт относится к многостраничным.

\newpage
\section{Тестирование совместимости}
% % --- browserling.com.
% % При тестировании следует провести попарное
% % тестирование ОС+браузер, например ОС windows 10 и браузер Opera версии 104,
% % затем изменить браузер на Edge версии 118. 
% % Браузеры и версии браузеров выбираете самостоятельно. 
% % Следует провести 2 теста и сделать 2 скриншота.

% 1. BL: Version 147.0.7727.102 (Official Build) (64-bit) Chrome;  Windows 11 25H2 26200.8245
% 2. MY: Windows 10, Firefox 140

Для проведения кроссбраузерного тестирования использован специализированный онлайн-сервис \url{browserling.com}. Данный инструмент позволяет проверить \url{kzgroup.ru} на корректность отображения и функционирования веб-приложения в различных комбинациях операционных систем и браузеров.

\subsection*{Тест №1}
Операционная система \textit{Windows 10}, браузер \textit{Mozilla Firefox} версии 140. В ходе тестирования проверялась визуальная целостность интерфейса и доступность основных навигационных элементов. Результат зафиксирован на \cref{fig:bl1,fig:bl2}.

\begin{figure}[H]
    \centering
    \includegraphics[width=1\textwidth]{figs/bl1.png}
    \caption{Отображение веб-приложения в среде Windows 10, Firefox 140}
    \label{fig:bl1}
\end{figure}
\begin{figure}[H]
    \centering
    \includegraphics[width=1\textwidth]{figs/bl2.png}
    \caption{Отображение веб-приложения в среде Windows 10, Firefox 140}
    \label{fig:bl2}
\end{figure}

\subsection*{Тест №2}
Операционная система \textit{Windows 11}, браузер \textit{Google Chrome} версии 147. Результат представлен на \cref{fig:my1,fig:my2}.

\begin{figure}[H]
    \centering
    \includegraphics[width=0.6\textwidth]{figs/my1.png}
    \caption{Отображение веб-приложения в среде Windows 11, Chrome 147}
    \label{fig:my1}
\end{figure}
\begin{figure}[H]
    \centering
    \includegraphics[width=0.6\textwidth]{figs/my2.png}
    \caption{Отображение веб-приложения в среде Windows 11, Chrome 147}
    \label{fig:my2}
\end{figure}

\subsection*{Анализ}
Тестирование на одинаковых размерах устройств показало, что сайт сохраняет структурную целостность и работоспособность в выбранных конфигурациях.
Ошибок и багов верстки при смене браузерных движков и операционных систем не обнаружено.


\newpage
\section{Тестирование адаптивности}

Тестирование адаптивности веб-приложения проведено с использованием DevTools. Проверка осуществлялась для трех наиболее востребованных разрешений экрана: $1920 \times 1080$, $1366 \times 768$ и $360 \times 800$. 

Как видно на \cref{fig:1920+1080.png,fig:1366+768.png,fig:360+800.png}, элементы интерфейса корректно масштабируются в зависимости от размера окна просмотра. При переходе к разрешению $360 \times 800$ пропадает фоновое изображение кировского завода.

\begin{figure}[H]
    \centering
    \includegraphics[width=1\textwidth]{figs/1920+1080.png}
    \caption{$1920 \times 1080$}
    \label{fig:1920+1080.png}
\end{figure}
\begin{figure}[H]
    \centering
    \includegraphics[width=1\textwidth]{figs/1366+768.png}
    \caption{$1366 \times 768$}
    \label{fig:1366+768.png}
\end{figure}
\begin{figure}[H]
    \centering
    \includegraphics[width=1\textwidth]{figs/360+800.png}
    \caption{$360 \times 800$}
    \label{fig:360+800.png}
\end{figure}

\newpage
\section{Тестирование верстки}

Проверка валидности HTML-кода главной страницы выполнена с помощью сервиса W3 Markup Validation Service (\cref{fig:w3.png}). В ходе статического анализа выявлены предупреждения и ошибки. 

Среди проблем зафиксирована, к примеру, отсутствие обязательных атрибутов у некоторых тегов. Тем не менее, критических сбоев, которые препятствуют отображению страницы современными браузерами, не зафиксировано.

\begin{figure}[H]
    \centering
    \includegraphics[width=1\textwidth]{figs/w3.pdf}
    \caption{Тестирование верстки через \url{validator.w3.org}.}
    \label{fig:w3.png}
    \textit{Примечание: \cref{fig:w3.png} добавлен как pdf и может быть приближен для более детального рассмотрения.}
\end{figure}

\newpage
\section{Тестирование структуры}

Анализ структуры веб-приложения и проверка ссылочной целостности проведены с использованием онлайн-сервисов Siteliner и BuiltWith. 

Отчет Siteliner (\cref{fig:siteliner.pdf}) демонстрирует отсутствие неработающих ссылок на просканированных страницах. Это свидетельствует о корректной настройке внутренней навигации.

Инструментарий BuiltWith (\cref{fig:bw_analitics,fig:bw_cdn,fig:bw_cms,fig:bw_frameworks,fig:bw_js,fig:bw_ssl}) позволил детализировать внутреннюю структуру и подключенные модули. 

Выявлено использование систем веб-аналитики (Яндекс.Метрика, Facebook Pixel), сети доставки контента (jQuery CDN), системы управления контентом (1C-Bitrix), а также ряда JavaScript-библиотек (core-js, GSAP). Для обеспечения безопасности данных используется SSL-сертификат от GlobalSign. 

Суммарное число технология (\cref{fig:bw_techs}) равно 45.

\begin{figure}[H]
    \centering
    \includegraphics[width=1\textwidth]{figs/siteliner.pdf}
    \caption{SiteLiner}
    \label{fig:siteliner.pdf}
\end{figure}
\begin{figure}[H]
    \centering
    \includegraphics[width=1\textwidth]{figs/bw_analitics.png}
    \caption{BuiltWith Аналитика}
    \label{fig:bw_analitics}
\end{figure}
\begin{figure}[H]
    \centering
    \includegraphics[width=1\textwidth]{figs/bw_cdn.png}
    \caption{BuiltWith CDN}
    \label{fig:bw_cdn}
\end{figure}
\begin{figure}[H]
    \centering
    \includegraphics[width=1\textwidth]{figs/bw_cms.png}
    \caption{BuiltWith CMS}
    \label{fig:bw_cms}
\end{figure}
\begin{figure}[H]
    \centering
    \includegraphics[width=1\textwidth]{figs/bw_frameworks.png}
    \caption{BuiltWith FrameWorks}
    \label{fig:bw_frameworks}
\end{figure}
\begin{figure}[H]
    \centering
    \includegraphics[width=1\textwidth]{figs/bw_js.png}
    \caption{BuiltWith JavaScript}
    \label{fig:bw_js}
\end{figure}
\begin{figure}[H]
    \centering
    \includegraphics[width=1\textwidth]{figs/bw_ssl.png}
    \caption{BuiltWith SSL}
    \label{fig:bw_ssl}
\end{figure}
\begin{figure}[H]
    \centering
    \includegraphics[width=1\textwidth]{figs/bw_techs.png}
    \caption{BuiltWith Количество технология}
    \label{fig:bw_techs}
\end{figure}


\newpage
\section{Аудит веб-приложения}

Технический аудит выполнен с помощью профильного сервиса W3Techs (\cref{fig:w3techs.pdf}). 

Результаты сканирования подтверждают данные, полученные на предыдущем этапе: сайт работает на базе CMS Bitrix, в качестве серверного языка программирования применяется PHP 8.2.29, веб-сервером выступает Nginx и т. д.

Дополнительно зафиксировано использование кодировки UTF-8 и языка разметки HTML5.

\begin{figure}[H]
    \centering
    \includegraphics[width=1\textwidth]{figs/w3techs.pdf}
    \caption{W3Techs}
    \label{fig:w3techs.pdf}
    \textit{Примечание: \cref{fig:w3techs.pdf} добавлен как pdf и может быть приближен для более детального рассмотрения.}
\end{figure}

\newpage
\section{Выводы}

В ходе выполнения практической работы была достигнута поставленная цель: получены практические навыки тестирования клиентской части веб-приложения. 

В процессе работы были приобретены следующие навыки:
\begin{itemize}
    \item Проведение кроссбраузерного тестирования с использованием специализированных платформ для выявления дефектов отображения в различных операционных системах и версиях браузеров.
    \item Оценка адаптивности пользовательского интерфейса средствами DevTools при эмуляции различных разрешений экранов.
    \item Валидация HTML-разметки на соответствие стандартам W3 для обеспечения кроссплатформенной совместимости.
    \item Анализ целостности ссылок и проведение комплексного технического аудита стека технологий с применением Siteliner, BuiltWith и W3Techs.
\end{itemize}

В ходе тестирования возникла проблема, связанная с жестким ограничением времени сессии в 3 минуты на \url{browserling.com}. Данная проблема была решена путем предварительного планирования тестовых сценариев: конфигурации сред ОС + браузер были выбраны заранее, а создание скриншотов производилось немедленно после загрузки интерфейса. 

Итоговый анализ показал, что протестированное веб-приложение обладает высоким уровнем кроссбраузерной совместимости, адаптивности и технической оптимизации, не имеет нарушений ссылочной целостности, однако требует корректировки HTML-кода для устранения некоторых ошибок.